\section{\texorpdfstring{Coupling independence for Edge-Potts at \(q\ge3\Deg\)}
{Coupling independence for Edge-Potts at q >= 3 Delta}}
\label{sec:soft-edge}
\label{app:soft-edge-ci}

For a finite simple graph \(G=(V,E)\), edge-colour assignments are vertex
colourings of the line graph \(L(G)\).  An edge-colour pinning
\(\tau:\Lambda\to\colours\), with \(\Lambda\subseteq E\), is likewise a
vertex-colour pinning of \(L(G)\).  We write
\[
  (\edgegibbs{G}{x})^\tau:=\mupin{L(G)}{\tau}{x}.
\]
The Fisher polynomial of this edge-Potts model is simply the vertex-Potts
polynomial on \(L(G)\); in particular,
\begin{equation}\label{eq:edge-normalization}
  \Zpin{L(G)}{\tau}(z)
  =z^{m_{L(G)}(\tau)}\nZpin{L(G)}{\tau}(z).
\end{equation}
The analytic argument is the vertex-Potts transfer on line graphs; the
additional input is the following coupling bound.

\begin{theorem}[Edge-Potts coupling independence]
\label{thm:soft-edge-ci}
Let \(q,\Deg\) be positive integers with \(q\ge3\Deg\).  For every finite simple graph \(G\)
of maximum degree at most \(\Deg\), every edge-colour pinning \(\tau\),
every \(x\in[0,1]\), every free edge
\(i\in E\setminus\operatorname{dom}\tau\), and all
\(a,b\in\colours\), put, for \(c\in\colours\),
\[
  \tau^{i=c}:=\tau\cup\{i\mapsto c\},
  \qquad
  (\edgegibbs{G}{x})^{\tau,i=c}
  :=(\edgegibbs{G}{x})^{\tau^{i=c}}.
\]
Then, on the common free edge set \(E\setminus(\operatorname{dom}\tau\cup\{i\})\),
\begin{equation}\label{eq:soft-edge-ci}
  \WHam\!\left(
    (\edgegibbs{G}{x})^{\tau,i=a},
    (\edgegibbs{G}{x})^{\tau,i=b}
  \right)
  \le\Deg-1.
\end{equation}
\end{theorem}

\begin{corollary}[Edge-Potts zero-freeness]
\label{cor:soft-edge-zf}
Let \(q,\Deg\) be positive integers with \(\Deg\ge2\) and
\(q\ge3\Deg\).  There exists
\(\eps=\eps(q,\Deg)>0\) such that, for every finite simple graph \(G\) of
maximum degree at most \(\Deg\) and every edge-colour pinning \(\tau\),
\[
  \nZpin{L(G)}{\tau}(z)\ne0
  \qquad\text{whenever }\dist(z,[0,1])<\eps.
\]
Thus the unnormalized pinned polynomial has no zero in this neighbourhood
except possibly at \(z=0\); its order of vanishing there is
\(m_{L(G)}(\tau)\).  In particular,
\(Z_{L(G)}(z)\ne0\) throughout the same neighbourhood.
\end{corollary}

\begin{proof}
The line graph has maximum degree at most
\(\Deg_L:=2\Deg-2\).  The class of line graphs of simple graphs of maximum
degree at most \(\Deg\) is closed under taking induced subgraphs, since an
induced subgraph on an edge set \(S\) is the line graph of \((V,S)\).
Vertex pinnings in this class are exactly edge-colour pinnings, with the
normalization in \eqref{eq:edge-normalization}.
Moreover, \(q\ge3\Deg\) implies \(q\ge\Deg_L+1\), and
\Cref{thm:soft-edge-ci} supplies one coupling-independence constant on
\([0,1]\).  Apply
\Cref{thm:potts-transfer} with degree bound \(\Deg_L\).  The statement for
the unnormalized polynomial follows from \eqref{eq:edge-normalization}.
\end{proof}

\paragraph{Proof overview.}

For \(0<x<1\), hidden endpoint slots represent the edge-Potts weight by a
hard distinct-label constraint.  Replacing one occupied endpoint label
admits a coupling with expected Hamming distance at most \((\Deg-1)/2\);
conditioning \(i\) requires one replacement at each endpoint.  The resulting
coupling is a countable weighted adaptation of the list-edge-colouring
coupling in~\cite{CWZZ2025}.

\subsection{Exact slot lift}

Fix a finite simple graph \(G=(V,E)\) of maximum degree at most
\(\Deg\), and \(q\ge3\Deg\).  Fix \(0<x<1\), a subset
\(\Lambda\subseteq E\), and an edge-colour pinning
\(\tau\colon\Lambda\to\colours\).  Thus
\(\operatorname{dom}\tau=\Lambda\), and
\(F_0=E\setminus\Lambda\) is the initial set of free edges.  Fix once and for all an
ordering \(e=(e^-,e^+)\) of the two endpoints of every \(e\in E\).  When
using slot notation for \(e\), we always write
\((u,v)=(e^-,e^+)\), so that \(e=uv\) has this fixed ordered meaning.
For a vertex
\(v\) and colour \(c\), write
\[
  \bcount_v^\tau(c)
  :=\abs{\set{e\in\operatorname{dom}\tau:e\ni v,\ \tau(e)=c}},
\]
so here $\bcount_v^\tau(c)$ counts pinned edges incident to the base-graph
vertex $v$.  For \(\phi\colon F_0\to\colours\), put
\[
  k_{v,c}(\phi)
  :=\abs{\set{e\in F_0:e\ni v,\ \phi(e)=c}}.
\]
For an edge colouring \(\psi:D\to\colours\), where \(D\subseteq E\), put
\[
  \Nconf(\psi)
  :=\abs{\set{\{e,f\}\subseteq D:
        e\ne f,\ e\cap f\ne\varnothing,\ \psi(e)=\psi(f)}}.
\]
Then \(m_{L(G)}(\tau)=\Nconf(\tau)\).
Because \(G\) is simple, two distinct edges share at most one endpoint.
Consequently,
\begin{equation}\label{eq:edge-conflict-split}
  \Nconf(\tau\cup\phi)-m_{L(G)}(\tau)
  =
  \sum_{v,c}
  \left[
    \bcount_v^\tau(c)\,k_{v,c}(\phi)
    +\binom{k_{v,c}(\phi)}2
  \right].
\end{equation}
The two terms count pinned--free and free--free conflicts, respectively;
the pinned--pinned conflicts are precisely the subtracted normalization.

The second term in \eqref{eq:edge-conflict-split} has the following
probabilistic representation.

\begin{lemma}[Slot representation]\label{lem:slot-fact}
For every \(0<x<1\), there is a probability distribution
\((\kappa_r)_{r\ge1}\), with every \(\kappa_r>0\), such that for every integer
\(k\ge0\),
\begin{equation}\label{eq:slot-identity}
  \sum_{\substack{r_1,\ldots,r_k\ge1\\
                  r_1,\ldots,r_k\ {\rm all\ distinct}}}
       \prod_{j=1}^k \kappa_{r_j}
  =x^{\binom k2}.
\end{equation}
\end{lemma}

\begin{proof}
Define the deformed exponential

\[
  \mathcal E_x(t):=\sum_{k\ge0}x^{\binom k2}\frac{t^k}{k!}.
\]

By~\cite{MFB1972,Iserles1993}, all zeros of \(\mathcal E_x\) are real and
negative; see also~\cite[\S1, pp.~312--313]{WangZhang2018}.  Write them,
with multiplicity, as \(-1/\kappa_r\), where \(\kappa_r>0\).  If \(a_k\)
is the coefficient of \(t^k\), then

\[
  \log\frac1{\abs{a_k}}
  =\abs{\log x}\binom k2+\log k!,
\]

so \(\mathcal E_x\) is entire of order zero.  Its genus-zero Hadamard
factorization is

\[
  \mathcal E_x(t)
  =\prod_{r\ge1}(1+\kappa_r t).
\]

There is no exponential factor, and the constant is fixed by
\(\mathcal E_x(0)=1\).  In particular,
\(\sum_r\kappa_r<\infty\), so the product and its coefficient
expansions converge absolutely.
Comparison of the linear coefficients gives \(\sum_r\kappa_r=1\), and
comparison at degree \(k\) gives

\[
  \sum_{1\le r_1<\cdots<r_k}
       \kappa_{r_1}\cdots\kappa_{r_k}
  =\frac{x^{\binom k2}}{k!}.
\]

Multiplication by \(k!\) gives \eqref{eq:slot-identity}.
\end{proof}

Thus \(k\) independent slots with law \((\kappa_r)\) are all distinct with
probability \(x^{\binom k2}\).  Apply this identity separately at every
vertex--colour pair.

For a free edge \(e=uv\), use the lifted state
\[
  \xi=(c,r,s)\in\colours\times\mathbb N_{\ge1}^2
\]
with activity
\begin{equation}\label{eq:edge-slot-activity}
  w_e(c,r,s)
  :=x^{\bcount_u^\tau(c)+\bcount_v^\tau(c)}\kappa_r\kappa_s.
\end{equation}
These activities always use the original colour pinning $\tau$ and remain
fixed during later lifted-state pinnings; the latter are recorded only in
the occupied-label sets defined below.
At \(u\) this state exposes the endpoint label
\(\ell_{e,u}(\xi)=(c,r)\), and at \(v\) it exposes
\(\ell_{e,v}(\xi)=(c,s)\).  A lifted assignment is admissible when all
labels exposed at any one vertex are distinct.

After some lifted edge states have been pinned, let \(F\subseteq F_0\)
denote the current free-edge set and let \(B_v\subseteq
\colours\times\mathbb N_{\ge1}\) be a finite \emph{set} of occupied labels
at \(v\), write \(B=(B_v)_{v\in V}\), and put
\[
  \deg_F(v):=\abs{\set{e\in F:v\in e}}.
\]
An assignment
\[
  \omega=(\omega_e)_{e\in F}
  \in\bigl(\colours\times\mathbb N_{\ge1}^2\bigr)^F
\]
on the current free-edge set is admissible when
the labels in \(B_v\), together with those exposed by the incident free
edges, are pairwise distinct at every \(v\).  Its weight is
\[
  W_{F,B}(\omega):=\prod_{e\in F}w_e(\omega_e),
\]
and \(\nu_{F,B}\) denotes the corresponding normalized law on admissible
assignments whenever its partition sum is positive.

For a current free edge \(e=uv\), define
\[
  d_F(e)
  :=\abs{\set{f\in F\setminus\set e:f\cap e\ne\varnothing}}
\]
and let
\[
  M_e(B)
  :=
  \sum_{\substack{\xi=(c,r,s)\\
       \ell_{e,u}(\xi)\notin B_u,\;
       \ell_{e,v}(\xi)\notin B_v}}
       w_e(\xi)
\]
be the activity compatible with the occupied boundary labels.  The
required bounds, for every \(v\in V\) and \(e\in F\) as applicable, are
\begin{align}
  \deg_F(v)+\abs{B_v}&\le\Deg,
  \label{eq:slot-load}\\
  \sum_{\xi:\,\ell_{e,y}(\xi)=\zeta}w_e(\xi)&\le1
  \quad
  (y\in e,\ \zeta\in\colours\times\mathbb N_{\ge1}),
  \label{eq:slot-fibre}\\
  M_e(B)-d_F(e)&\ge\Deg+2.
  \label{eq:slot-slack}
\end{align}
Pinning an admissible positive-activity state \(\xi\) on \(e=uv\) produces
\(F'=F\setminus\set e\) and the occupied-label family \(B'\) given by
\[
 B'_u=B_u\cup\set{\ell_{e,u}(\xi)},\qquad
 B'_v=B_v\cup\set{\ell_{e,v}(\xi)},\qquad
 B'_y=B_y\quad(y\notin\{u,v\}).
\]
\begin{lemma}[Exact lift and inherited bounds]\label{lem:edge-slot-lift}
Assume \(q\ge3\Deg\).  For the initial data
\((F_0,\varnothing)\), where $\varnothing$ denotes the family of empty
occupied-label sets, colour projection
\((c,r,s)\mapsto c\) sends the lifted law exactly to the normalized pinned
edge-Potts law \((\edgegibbs{G}{x})^\tau\).

The initial slot instance and every boundary obtained by successively
pinning admissible positive-activity states satisfy
\eqref{eq:slot-load}--\eqref{eq:slot-slack}.  More generally, any instance
\((F,B)\) with finite occupied-label sets that uses the endpoint-label maps
and activities in \eqref{eq:edge-slot-activity} and satisfies these three
bounds has a finite, positive partition sum.
\end{lemma}

\begin{proof}
Fix a colour assignment \(\phi\colon F_0\to\colours\) and sum over its slot
fibre.  The sum factorizes over the pairs \((v,c)\).  If
\(k=k_{v,c}(\phi)\), the activities contribute
\(x^{\bcount_v^\tau(c)k}\), while admissibility requires the \(k\) endpoint
slots to be distinct.  By \eqref{eq:slot-identity}, their total slot weight
is \(x^{\binom k2}\).  Hence the fibre weight is
\[
  \prod_{v,c}
    x^{\bcount_v^\tau(c)k_{v,c}(\phi)+
       \binom{k_{v,c}(\phi)}2}
  =
  x^{\Nconf(\tau\cup\phi)-m_{L(G)}(\tau)}
\]
by \eqref{eq:edge-conflict-split}.  This proves exact projection.

Initially, \eqref{eq:slot-load} follows from
\(\deg_{F_0}(v)\le\Deg\).  For a fixed endpoint label, say \((c,r)\) at \(u\),
\[
  \sum_s w_e(c,r,s)
  =
  x^{\bcount_u^\tau(c)+\bcount_v^\tau(c)}\kappa_r\sum_s \kappa_s
  \le1,
\]
and the other endpoint is symmetric.  This proves
\eqref{eq:slot-fibre}.

For the slack, let
\[
  \bcount_e:=\sum_c\bigl(\bcount_u^\tau(c)+\bcount_v^\tau(c)\bigr),
\]
the number of pinned edges adjacent to \(e=uv\).  Since
\(\sum_r\kappa_r=1\),
\[
  M_e(\varnothing)
  =\sum_c x^{\bcount_u^\tau(c)+\bcount_v^\tau(c)}
  \ge q-\bcount_e.
\]
Here we used \(x^m\ge1-m\) for nonnegative integers \(m\) and \(0<x<1\).
Simplicity gives
\[
  d_{F_0}(e)+\bcount_e
  =
  \abs{\set{f\in E\setminus\set e:f\cap e\ne\varnothing}}
  \le2\Deg-2,
\]
and therefore
\[
  M_e(\varnothing)-d_{F_0}(e)
  \ge q-d_{F_0}(e)-\bcount_e
  \ge3\Deg-(2\Deg-2)=\Deg+2.
\]

It remains to check one pinning step for current data \((F,B)\).
Suppose \(e\) is pinned and a
remaining edge \(f\) meets \(e\).  Simplicity makes the common endpoint
unique.  The new occupied label removes from \(M_f\) at most one endpoint
fibre, whose activity is at most one by \eqref{eq:slot-fibre}; deleting
\(e\) decreases \(d_F(f)\) by exactly one.  Thus
\[
  M_f(B')-d_{F\setminus\set e}(f)
  \ge (M_f(B)-1)-(d_F(f)-1)
  \ge\Deg+2.
\]
Nonadjacent edges are unchanged.  At each endpoint of \(e\), the free
degree drops by one and one genuinely new occupied label is added, so
\eqref{eq:slot-load} is preserved.  The activities, and hence
\eqref{eq:slot-fibre}, do not change.

Finally, the three bounds imply feasibility by greedy exposure.  At every
step an unassigned edge has at least
\(M_e(B)-d_F(e)\ge\Deg+2>0\) compatible activity, and the preceding
pinning argument preserves the same conclusion.  The partition sum is
finite because there are finitely many edges and
\(\sum_\xi w_e(\xi)\le q\) for each edge.
\end{proof}

\subsection{Coupling under a boundary-label replacement}

Consider two occupied-label systems with the same free graph,
endpoint-label maps, and edge activities.  Fix a vertex \(v\), a finite
\(C_v\subseteq\colours\times\mathbb N_{\ge1}\) and distinct
\(\alpha,\beta\in
(\colours\times\mathbb N_{\ge1})\setminus C_v\).
Their occupied-label families satisfy
\(B_y^\alpha=B_y^\beta\) for \(y\ne v\), while at \(v\),
\[
  B_v^\alpha=C_v\cup\set\alpha,
  \qquad
  B_v^\beta=C_v\cup\set\beta.
\]
Both systems satisfy \eqref{eq:slot-load}--\eqref{eq:slot-slack}.  These
three bounds, rather than derivation from the initial lift, are the
hypotheses because the two-step comparison uses mixed endpoint boundaries.

For two lifted assignments on \(F\), let
\[
  \Ham_F(\omega,\omega')
  :=\bigl|\{e\in F:\omega_e\ne\omega'_e\}\bigr|,
\]
and write \(W_{1,F}\) for the infimum of the expected \(\Ham_F\)-cost
over couplings of two laws on the resulting countable state space.

For each integer \(n\ge0\), let \(\mathsf T_n\) be the supremum of
\(W_{1,F}(\nu_{F,B^\alpha},\nu_{F,B^\beta})\) over all data described in
the preceding paragraph that satisfy the three displayed bounds and have
at most \(n\) free edges.  The laws exist by \Cref{lem:edge-slot-lift},
\(0\le\mathsf T_n\le n\), and \(\mathsf T_0=0\).

\begin{lemma}[One-label coupling bound]\label{lem:edge-one-label}
One has \(\mathsf T_0=0\), and for every integer \(n\ge1\),
\begin{equation}\label{eq:edge-one-label-bound}
  \mathsf T_n
  \le
  \frac{\Deg-1}{2}
  \left[
    1-\left(\frac{\Deg-1}{\Deg}\right)^n
  \right]
  \le\frac{\Deg-1}{2}.
\end{equation}
\end{lemma}

\begin{proof}
Only \(n\ge1\) requires proof.  Fix an admissible comparison with
$1\le |F|\le n$; the case $F=\varnothing$ has zero cost.  Write
\(\mu_\alpha:=\nu_{F,B^\alpha}\),
\(\mu_\beta:=\nu_{F,B^\beta}\), and let
\(N=\set{j\in F:j\ni v}\).  Under \(\mu_\alpha\), define
\[
  E_j^\beta
  :=\set{\omega:\ell_{j,v}(\omega_j)=\beta},
  \qquad
  E_0^\beta
  :=\set{\omega:\text{no edge of \(N\) exposes \(\beta\) at \(v\)}}.
\]
Define \(E_j^\alpha,E_0^\alpha\) symmetrically under \(\mu_\beta\).
Distinctness of endpoint labels makes the label-exposure events
\((E_j^\beta)_{j\in N}\) disjoint, and together with \(E_0^\beta\) they
partition the \(\alpha\)-side state space; the same holds on the other
side.

The avoidance events have positive probability.  Fix either current
boundary \(B\in\{B^\alpha,B^\beta\}\).  For that side, greedily expose
the free edges while additionally forbidding the missing label at \(v\).
For an incident edge this removes at most one further endpoint fibre, so
the available activity is at least
\[
  M_e(B)-d_F(e)-1\ge\Deg+1>0.
\]
At each exposure, deleting the chosen edge lowers the remaining free
edge-degrees by the same amount as the new occupied endpoint labels can
lower their allowed activities, so this positive margin persists at every
later step.
Set
\[
  \rho_j
  :=\frac{\mu_\alpha(E_j^\beta)}{\mu_\alpha(E_0^\beta)},
  \qquad
  \sigma_j
  :=\frac{\mu_\beta(E_j^\alpha)}{\mu_\beta(E_0^\alpha)},
\]
with a zero numerator interpreted as zero.

To bound these label-exposure odds, start from a configuration in
\(E_j^\beta\), keep every state off \(j\) and replace the state of \(j\)
by a compatible state that does not expose \(\beta\) at \(v\).  The other
adjacent free edges exclude at most \(d_F(j)\) endpoint fibres, and
avoiding \(\beta\) excludes one more.  Thus the replacement activity is at
least \(\Deg+1\).  Route the source weight among the replacements
proportionally to their activities.  Conversely, for a fixed target, summing
all source states whose endpoint colour is \(\beta\) gives total activity
\(x^{\bcount_u^\tau(\beta)+\bcount_v^\tau(\beta)}\le1\),
which is also the bound supplied by \eqref{eq:slot-fibre}.  Summing this flow gives
\[
  \sum_{\omega\in E_j^\beta}W_{F,B^\alpha}(\omega)
  \le\frac1{\Deg+1}
       \sum_{\omega\in E_0^\beta}W_{F,B^\alpha}(\omega),
\]
and hence
\[
  \rho_j\le\frac1{\Deg+1}.
\]
The same argument gives \(\sigma_j\le1/(\Deg+1)\).  Since one boundary
label is already occupied at \(v\), \eqref{eq:slot-load} implies
\(\abs N\le\Deg-1\).  Hence
\begin{equation}\label{eq:edge-occurrence-budget}
  R:=\sum_{j\in N}\max\{\rho_j,\sigma_j\}
  \le\frac{\Deg-1}{\Deg+1}.
\end{equation}

Write
\[
  \nu_0^\alpha:=\mu_\alpha(\,\cdot\mid E_0^\beta),
  \qquad
  \nu_j^\alpha:=\mu_\alpha(\,\cdot\mid E_j^\beta),
\]
and define \(\nu_0^\beta,\nu_j^\beta\) symmetrically.  A conditional law
with zero coefficient below is omitted.  We next record the three
conditional comparisons:
\begin{align}
  W_{1,F}(\nu_0^\alpha,\nu_0^\beta)&=0,
  \label{eq:edge-avoidance-cost}\\
  W_{1,F}(\nu_j^\alpha,\nu_j^\beta)&\le1+\mathsf T_{n-1},
  \label{eq:edge-matched-cost}\\
  \max\!\left\{
    W_{1,F}(\nu_j^\alpha,\mu_\beta),
    W_{1,F}(\mu_\alpha,\nu_j^\beta)
  \right\}&\le1+2\mathsf T_{n-1}.
  \label{eq:edge-unmatched-cost}
\end{align}
For \eqref{eq:edge-avoidance-cost}, the two avoidance events produce the
same boundary \(C_v\cup\set{\alpha,\beta}\), so the conditional laws are
identical.

For the other two bounds, condition further on the complete state of
\(j=vw\), remove \(j\), and put its endpoint labels into the boundary.
For each pair of conditioned states, the remaining-edge laws are compared
through the corresponding residual boundary; mixing these conditional
couplings over the state of $j$ preserves the stated uniform cost.
For \eqref{eq:edge-matched-cost}, the boundaries at \(v\) then agree, and
the two systems differ by at most one label replacement at \(w\).  The
residual cost is at most \(\mathsf T_{n-1}\), and restoring \(j\) costs at
most one.

For \eqref{eq:edge-unmatched-cost}, compare the two residual laws through
the boundary that uses the first conditioned system's labels at \(v\) and
the second conditioned system's labels at \(w\).  This requires at most
one label replacement at each endpoint.  After \(j\) is removed, simplicity ensures
that no remaining edge meets both \(v\) and \(w\).  Every remaining edge
therefore sees exactly the endpoint restrictions of one of the two
conditioned systems, so the intermediate boundary satisfies
\eqref{eq:slot-load}--\eqref{eq:slot-slack} and is feasible.  Its residual
cost is at most \(2\mathsf T_{n-1}\), and restoring \(j\) again costs at
most one.
The conditioning mixtures are countable, but all these costs are uniform.
For each \(\eta>0\), choose every conditional coupling with expected cost
within \(\eta\) of its Wasserstein infimum.  Tonelli's theorem, followed by
\(\eta\downarrow0\), gives
\eqref{eq:edge-matched-cost}--\eqref{eq:edge-unmatched-cost}.

It remains to couple the exposure blocks.  The following table specifies
a nonnegative mixture of couplings.  Every displayed mass is divided by
\(1+R\); the first row is used once and the other rows once for each
$j\in N$.  Zero-mass rows are omitted, and $(t)_+:=\max\{t,0\}$:
\[
\begin{array}{c|c|c}
 \text{numerator}&\alpha\text{-side}&\beta\text{-side}\\ \hline
 1&\nu_0^\alpha&\nu_0^\beta\\
 \min\{\rho_j,\sigma_j\}&\nu_j^\alpha&\nu_j^\beta\\
 (\rho_j-\sigma_j)_+&\nu_j^\alpha&\mu_\beta\\
 (\sigma_j-\rho_j)_+&\mu_\alpha&\nu_j^\beta
\end{array}
\qquad(j\in N).
\]
To check the first marginal, note that
\[
  \nu_0^\alpha+\sum_j\rho_j\nu_j^\alpha
  =\left(1+\sum_j\rho_j\right)\mu_\alpha
\]
and
\[
  R=\sum_j\rho_j+\sum_j(\sigma_j-\rho_j)_+.
\]
Thus the rows on the \(\alpha\)-side sum to
\((1+R)\mu_\alpha\); the other marginal is symmetric.  Combining the table
with \eqref{eq:edge-avoidance-cost}--\eqref{eq:edge-unmatched-cost} gives
\[
\begin{aligned}
  W_{1,F}(\mu_\alpha,\mu_\beta)
  &\le
  \frac{
    \sum_j\min\{\rho_j,\sigma_j\}(1+\mathsf T_{n-1})
    +\sum_j\abs{\rho_j-\sigma_j}(1+2\mathsf T_{n-1})
  }{1+R}\\
  &\le
  \frac{R}{1+R}(1+2\mathsf T_{n-1})
  \le
  \frac{\Deg-1}{2\Deg}(1+2\mathsf T_{n-1}),
\end{aligned}
\]
where the last step uses \eqref{eq:edge-occurrence-budget}.  Taking the
supremum over all such comparisons gives
\[
 \mathsf T_n\le\frac{\Deg-1}{2\Deg}(1+2\mathsf T_{n-1}).
\]
Solving this affine recursion from \(\mathsf T_0=0\) yields
\eqref{eq:edge-one-label-bound}.
\end{proof}

\subsection{From slot couplings to coupling independence}

\begin{proof}[Proof of \Cref{thm:soft-edge-ci}]
If \(a=b\), the two normalized children are identical.  Assume
\(a\ne b\), and write the root edge as \(i=uv\).

\emph{Case \(x=1\).}
Under every pinning, the remaining edge colours are independent and
uniform, so the root-excluded distance is zero.

\emph{Case \(0<x<1\).}
Construct the single slot lift associated with the original pinning
\(\tau\).  By \Cref{lem:edge-slot-lift}, its colour projection is exactly
\((\edgegibbs{G}{x})^\tau\).  Conditioning the root edge \(i\) to
colour \(a\) disintegrates the lifted child over the positive-mass root
states \((a,r,s)\); the \(b\)-child is the analogous mixture over
\((b,r',s')\).  Both are conditionals of this same lift.

For \(c\in\{a,b\}\), let \(\Pi_c\) denote the law of \(\omega_i\) under
\(\nu_{F_0,\varnothing}\), conditional on the colour coordinate of
\(\omega_i\) being \(c\).  Thus every
\(\xi\in\operatorname{supp}\Pi_c\) has the form \((c,r,s)\).  For such a
state \(\xi\), delete \(i\) and define the occupied-label family on the
remaining system by
\[
  B_y^\xi:=
  \begin{cases}
    \set{\ell_{i,y}(\xi)},&y\in i,\\
    \varnothing,&y\notin i.
  \end{cases}
\]
Set
\[
  \nu^\xi:=\nu_{F_0\setminus\{i\},B^\xi}.
\]
Thus the root-excluded lifted \(c\)-child is the \(\Pi_c\)-mixture of the
laws \(\nu^\xi\).

Fix \(\xi\in\operatorname{supp}\Pi_a\) and
\(\xi'\in\operatorname{supp}\Pi_b\), and delete \(i\).  Each conditioned
law satisfies
\eqref{eq:slot-load}--\eqref{eq:slot-slack}.  Compare the two laws through
the boundary that takes the first root state's label at \(u\) and the
second root state's label at \(v\).  Since \(G\) is simple, no remaining
edge meets both \(u\) and \(v\); hence this intermediate boundary inherits
the three bounds from the two conditioned systems.  The two comparisons
change at most one occupied label each.  Therefore
\[
  W_{1,F_0\setminus\{i\}}(\nu^\xi,\nu^{\xi'})
  \le2\mathsf T_{\abs{F_0}-1}\le\Deg-1
\]
by \Cref{lem:edge-one-label}.

Couple the two countable root-state laws by the product law
\(\Pi_a\otimes\Pi_b\).  For each \(\eta>0\) and every positive-mass pair
\((\xi,\xi')\), choose a coupling of \(\nu^\xi\) and \(\nu^{\xi'}\) whose
expected cost is within \(\eta\) of the corresponding Wasserstein infimum,
and then let \(\eta\downarrow0\).  Tonelli's theorem preserves the same
uniform bound.  Colour projection cannot increase Hamming distance, so the
two edge-Potts children have root-excluded distance at most \(\Deg-1\).

\emph{Case \(x=0\).}
Work in the finite colour space.  Every normalized hard child exists:
after fixing \(i\), greedily colour the remaining vertices of
the line graph, equivalently the remaining free edges of \(G\).  Each sees
at most \(2\Deg-2\) already coloured neighbours,
whereas \(q\ge3\Deg\).

For \(c\in\set{a,b}\), put
\[
  t_c:=\bcount_u^\tau(c)+\bcount_v^\tau(c).
\]
Then
\[
  m_{L(G)}(\tau^{i=c})
  =m_{L(G)}(\tau)+t_c,
\]
and for every colouring \(\phi\) of the remaining edges and \(x>0\),
\[
\begin{aligned}
  &x^{\Nconf(\tau^{i=c}\cup\phi)
       -m_{L(G)}(\tau)}\\
  &\qquad
  =
  x^{t_c}\,
  x^{\Nconf(\tau^{i=c}\cup\phi)
       -m_{L(G)}(\tau^{i=c})}.
\end{aligned}
\]
The first factor is independent of \(\phi\), so ordinary conditioning for
\(x>0\) is exactly the normalized child.  After this scalar is removed,
the two families of weights are finite polynomials in \(x\), and their
partition sums are positive at \(x=0\) by the preceding greedy argument.
\Cref{lem:finite-endpoint-closure} therefore passes the bound
\(\Deg-1\) to the normalized hard children.
\end{proof}
